\documentclass[11pt]{article}

\usepackage[margin=1in]{geometry}
\usepackage{amsmath,amssymb,amsthm}
\usepackage{algorithm}
\usepackage{algpseudocode}
\usepackage{enumitem}
\usepackage{hyperref}
\usepackage{xcolor}
\usepackage{tikz}
\usepackage{booktabs}

\hypersetup{colorlinks=true, linkcolor=blue, urlcolor=blue, citecolor=blue}

\newtheorem{theorem}{Theorem}[section]
\newtheorem{lemma}[theorem]{Lemma}
\newtheorem{proposition}[theorem]{Proposition}
\newtheorem{corollary}[theorem]{Corollary}
\newtheorem{definition}[theorem]{Definition}
\newtheorem{example}[theorem]{Example}

\title{TOAE209/TOAH209 -- Design and Analysis of Algorithms\\[4pt]
\large Week 11: Network Flow II/III: Bipartite Matching and Applications}
\author{Foreign Trade University}
\date{}

\begin{document}
\maketitle
\tableofcontents

\section{Introduction: Flow as a Modeling Language}

Last week we developed the theory of \textbf{maximum flow} and the
\textbf{max-flow min-cut theorem}: in any flow network the value of a maximum
$s$--$t$ flow equals the capacity of a minimum $s$--$t$ cut, and the
Ford--Fulkerson family of algorithms (in particular Edmonds--Karp, which always
augments along a shortest path) computes both. That theorem is powerful in its
own right, but its real importance is as a \emph{modeling language}: a
surprising number of combinatorial problems that have nothing to do with
``pipes and fluid'' can be \emph{reduced} to a max-flow or a min-cut
computation. This week (Kleinberg \& Tardos, Ch.\ 7) is about those reductions.

The recurring pattern is always the same:
\begin{enumerate}[label=(\arabic*)]
    \item Start with a combinatorial problem (a matching, a set of disjoint
    paths, a selection of projects, an image to segment).
    \item \textbf{Build} a flow network whose structure encodes the constraints
    of the problem.
    \item \textbf{Prove} a correspondence: feasible/optimal solutions of the
    original problem correspond exactly to flows/cuts of a certain value.
    \item \textbf{Solve} the flow problem and \textbf{translate} the answer
    back.
\end{enumerate}

The art is entirely in steps (2) and (3); step (4) is a black-box call to last
week's algorithm. Throughout, we use the shared \texttt{MaxFlow} helper (an
Edmonds--Karp implementation in \texttt{starter\_code.py}) as that black box, so
the practical problems can focus on the \emph{reductions}.

\subsection{Roadmap}

We begin with the cleanest and most important application -- \textbf{bipartite
matching} -- and the two beautiful min-max theorems that surround it
(\textbf{Hall's theorem} and \textbf{K\H{o}nig's theorem}). We then treat
\textbf{disjoint paths} (Menger's theorem), \textbf{circulations with demands
and lower bounds} (a more expressive flow model), and finally a sequence of
``killer app'' reductions: \textbf{bipartite assignment / survey design},
\textbf{baseball elimination}, \textbf{project selection} (max-weight closure
via min-cut), and \textbf{image segmentation} (also a min-cut). Practical
Problems 1--12 implement essentially all of these.

\section{Bipartite Matching via Max Flow}

\subsection{Problem statement}

\begin{definition}[Bipartite graph]
A graph $G=(V,E)$ is \emph{bipartite} if $V$ can be partitioned into two sets
$X$ and $Y$ such that every edge has one endpoint in $X$ and one in $Y$.
Equivalently (and this is what Problem 1 checks), $G$ is bipartite iff it has no
odd-length cycle, iff its vertices can be 2-colored so adjacent vertices differ.
\end{definition}

\begin{definition}[Matching]
A \emph{matching} $M \subseteq E$ is a set of edges no two of which share an
endpoint. A vertex is \emph{matched} (or \emph{saturated}) if it is an endpoint
of some edge of $M$, and \emph{free} otherwise. A \textbf{maximum matching} is a
matching of largest possible size. A \emph{perfect matching} saturates every
vertex.
\end{definition}

The \textbf{Maximum Bipartite Matching Problem} asks for a maximum matching in a
given bipartite graph $G=(X \cup Y, E)$. This models any ``pairing'' task:
assigning workers to jobs, students to projects, applicants to positions, etc.

\subsection{The flow reduction}

\begin{algorithm}[h]
\caption{Bipartite Matching via Max Flow}
\begin{algorithmic}[1]
\State Create a source $s$ and sink $t$.
\State For each $x \in X$: add edge $s \to x$ with capacity $1$.
\State For each $y \in Y$: add edge $y \to t$ with capacity $1$.
\State For each edge $(x,y) \in E$: add edge $x \to y$ with capacity $1$.
\State Compute a maximum integer $s$--$t$ flow $f$.
\State \Return the set of edges $(x,y) \in E$ carrying one unit of flow.
\end{algorithmic}
\end{algorithm}

The capacity-$1$ edges out of $s$ and into $t$ enforce that each $x$ is matched
at most once and each $y$ is matched at most once. Problem 3 implements exactly
this reduction; Problem 2 implements the direct combinatorial alternative
(Kuhn's augmenting-path algorithm), and Problem 3's tests \emph{verify the two
give equal matching size on random graphs}.

\begin{theorem}[Matching = integer flow]
\label{thm:matching-flow}
The maximum matching size in $G$ equals the value of a maximum $s$--$t$ flow in
the network above. Moreover, every integer flow corresponds to a matching, and
vice versa.
\end{theorem}

\begin{proof}
\textbf{Matching $\Rightarrow$ flow.} Given a matching $M$ of size $k$, route one
unit of flow along $s \to x \to y \to t$ for each $(x,y) \in M$. Since $M$ is a
matching, no two of these paths share an $x$ or a $y$, so every capacity-$1$
edge carries at most one unit; this is a feasible flow of value $k$.

\textbf{Flow $\Rightarrow$ matching.} Because all capacities are integers, by the
integrality theorem (last week) there is an integer maximum flow $f$, in which
every edge carries $0$ or $1$ unit. The middle edges $x \to y$ carrying one unit
form a set $M$. Each $x$ has at most one unit entering it (from $s$) and hence at
most one unit leaving it, so $x$ is in at most one edge of $M$; likewise each
$y$. Thus $M$ is a matching, of size equal to the value of $f$.

The two directions show the maximum matching size equals the maximum flow value.
\end{proof}

The reduction runs in $O(V \cdot E)$ time with Edmonds--Karp (each augmenting
path increases the matching by one, and there are at most $\min(|X|,|Y|)$ of
them), matching the running time of Kuhn's algorithm. Problem 6 reuses this
machinery to answer an \emph{assignment feasibility} question: given people,
tasks, and an ``allowed'' relation, can \emph{everyone} be assigned a distinct
task? That is exactly asking whether a matching saturating all people exists.

\section{Hall's Theorem and K\H{o}nig's Theorem}

Two classical min-max theorems characterize bipartite matchings precisely. Both
can be derived from max-flow min-cut, and both are implemented this week
(Problems 4 and 5).

\subsection{Hall's marriage theorem}

\begin{definition}[Neighborhood and deficiency]
For a set $S \subseteq X$, its \emph{neighborhood} $N(S) \subseteq Y$ is the set
of all vertices adjacent to some vertex of $S$. We call $S$ \emph{deficient} if
$|N(S)| < |S|$.
\end{definition}

\begin{theorem}[Hall, 1935]
\label{thm:hall}
A bipartite graph $G=(X \cup Y, E)$ has a matching saturating every vertex of
$X$ (a \emph{perfect matching of $X$}) \textbf{if and only if} $|N(S)| \ge |S|$
for every subset $S \subseteq X$.
\end{theorem}

\begin{proof}
($\Rightarrow$) If a matching saturates $X$, then for any $S \subseteq X$ the
matched partners of the vertices of $S$ are $|S|$ distinct vertices, all lying in
$N(S)$; hence $|N(S)| \ge |S|$.

($\Leftarrow$) We prove the contrapositive using max-flow min-cut. Build the flow
network of Theorem~\ref{thm:matching-flow}. The matching saturates all of $X$
iff the max flow value equals $|X|$, iff (max-flow min-cut) the minimum cut has
capacity $|X|$. Consider any $s$--$t$ cut $(A,B)$ with $s \in A$. Let
$S = X \cap B$ (the $X$-vertices on the sink side) and note $X \cap A$
contributes capacity $|X \cap A|$ via the cut edges $s \to x$. To keep the cut
finite (the middle edges have capacity $1$ but we may instead cut the
$y \to t$ edges), an optimal finite cut places $N(S) \subseteq A$ on the source
side -- otherwise an infinite-capacity... more carefully: the cheapest cut that
``gives up'' on the vertices of $S$ must pay $1$ for each $s\to x$ with
$x \in X\cap A = X \setminus S$, namely $|X|-|S|$, plus $1$ for each $y \to t$
with $y \in N(S)$ on the source side, namely $|N(S)|$. So the minimum cut
capacity is $\min_{S \subseteq X}\big(|X| - |S| + |N(S)|\big)$. This is less than
$|X|$ exactly when some $S$ has $|N(S)| < |S|$, i.e.\ some deficient set exists.
Hence: no $X$-saturating matching $\iff$ a deficient set exists. Contrapositive:
if no set is deficient, an $X$-saturating matching exists.
\end{proof}

Problem 4 turns this proof into code: it decides whether $X$ has a perfect
matching and, when not, \emph{returns an explicit deficient set $S$} (recovered
from the alternating-path structure of a maximum matching). The deficient set is
a \emph{certificate} of infeasibility -- a short, checkable reason why no perfect
matching can exist.

\subsection{K\H{o}nig's theorem: matching = vertex cover}

\begin{definition}[Vertex cover]
A \emph{vertex cover} of $G=(V,E)$ is a set $C \subseteq V$ such that every edge
has at least one endpoint in $C$. The \textbf{minimum vertex cover} is a smallest
such set.
\end{definition}

For \emph{any} graph, a matching and a vertex cover are in tension:

\begin{lemma}[Weak duality]
\label{lem:weak-duality}
In any graph, $|M| \le |C|$ for every matching $M$ and every vertex cover $C$. In
particular, $\max$-matching $\le \min$-vertex-cover.
\end{lemma}

\begin{proof}
The edges of $M$ are pairwise disjoint, so they have $2|M|$ distinct endpoints.
Every edge of $M$ must be covered, i.e.\ have an endpoint in $C$; distinct edges
of $M$ need distinct cover vertices (they share no endpoint). Hence
$|C| \ge |M|$.
\end{proof}

For \emph{general} graphs this inequality can be strict: a triangle has maximum
matching $1$ but minimum vertex cover $2$ (Problem 12 exhibits exactly this). In
\emph{bipartite} graphs, equality always holds:

\begin{theorem}[K\H{o}nig, 1931]
\label{thm:konig}
In a bipartite graph, the size of a maximum matching equals the size of a
minimum vertex cover.
\end{theorem}

\begin{proof}
By Lemma~\ref{lem:weak-duality}, $\max$-matching $\le \min$-cover; we exhibit a
vertex cover of size exactly $|M|$, where $M$ is a maximum matching. Let
$U \subseteq X$ be the set of $X$-vertices left \emph{unmatched} by $M$. Mark
every vertex reachable from $U$ by an \emph{alternating path} (a path that
alternates non-matching edges $X \to Y$ and matching edges $Y \to X$). Let
$Z$ be the set of marked vertices, and define
\[
   C \;=\; (X \setminus Z) \,\cup\, (Y \cap Z).
\]
\emph{$C$ is a cover.} Suppose some edge $(x,y)$ has neither endpoint in $C$,
i.e.\ $x \in X \cap Z$ (marked) and $y \in Y \setminus Z$ (unmarked). If $(x,y)
\notin M$, then from the marked $x$ we could extend an alternating path along the
non-matching edge $(x,y)$ to reach $y$, marking it -- contradiction. If $(x,y)
\in M$, then $x$ was marked only by arriving along the matching edge $(x,y)$ from
$y$, so $y$ would be marked -- contradiction. Hence every edge is covered.

\emph{$|C| = |M|$.} Every vertex of $C$ is saturated by $M$: a vertex in
$X \setminus Z$ is matched (the unmatched $X$-vertices are all in $U \subseteq Z$),
and a vertex in $Y \cap Z$ is matched (an unmatched $Y$-vertex would complete an
augmenting path from $U$, contradicting maximality of $M$). Moreover no matching
edge has \emph{both} endpoints in $C$: if $(x,y) \in M$ had $x \in X\setminus Z$
and $y \in Y \cap Z$, then $y$ marked implies (since the only way to mark a
matched $y$ is via its matching partner) that $x$ is marked, i.e.\ $x \in Z$,
contradiction. So each edge of $M$ contributes exactly one endpoint to $C$, and
every vertex of $C$ is the endpoint of exactly one matching edge -- giving a
bijection between $M$ and $C$. Thus $|C| = |M|$.
\end{proof}

This proof is \emph{constructive}, and Problem 5 implements it directly: given a
maximum matching (from Problem 2), it builds a vertex cover of size $|M|$ via the
alternating-path marking, and the tests verify $|\text{cover}| = |\text{matching}|$
on random bipartite graphs. Problem 12 ties everything together by checking
$\max$-matching $\le \min$-cover for general graphs and equality for bipartite
graphs via brute force.

\begin{corollary}[Hall from K\H{o}nig]
Hall's theorem follows from K\H{o}nig's: $X$ has a perfect matching iff
$\min$-cover $= |X|$, and one checks that a cover smaller than $|X|$ exists iff
some $S \subseteq X$ is deficient.
\end{corollary}

\section{Disjoint Paths}

\subsection{Edge-disjoint paths and Menger's theorem}

\begin{definition}[Edge-disjoint paths]
A collection of $s$--$t$ paths is \emph{edge-disjoint} if no edge is used by more
than one path.
\end{definition}

\begin{theorem}[Menger, edge version]
\label{thm:menger-edge}
The maximum number of edge-disjoint $s$--$t$ paths in a directed graph equals the
minimum number of edges whose removal disconnects $t$ from $s$ (a minimum
$s$--$t$ edge cut).
\end{theorem}

\begin{proof}
Assign every edge capacity $1$ and compute a max flow. An integer flow of value
$k$ decomposes (by repeatedly tracing $s$--$t$ paths along positive-flow edges,
and discarding any cycles) into $k$ edge-disjoint $s$--$t$ paths, since each
capacity-$1$ edge can belong to at most one path. Conversely $k$ edge-disjoint
paths give a flow of value $k$. So $\max$ edge-disjoint paths $=$ max flow value
$=$ min cut capacity (max-flow min-cut) $=$ min number of edges to delete.
\end{proof}

Problem 7 implements this: build the unit-capacity network, return the max flow.

\subsection{Vertex-disjoint paths via node splitting}

To forbid two paths from sharing an \emph{internal vertex}, we use the
\textbf{node-splitting} gadget:

\begin{algorithm}[h]
\caption{Vertex-Disjoint Paths via Node Splitting}
\begin{algorithmic}[1]
\For{each internal vertex $v$ (i.e.\ $v \ne s,t$)}
    \State Split $v$ into $v_{\text{in}}$ and $v_{\text{out}}$, with an edge
    $v_{\text{in}} \to v_{\text{out}}$ of capacity $1$.
\EndFor
\State Replace each original edge $(u,v)$ by $u_{\text{out}} \to v_{\text{in}}$
of capacity $1$.
\State Compute a max $s$--$t$ flow; its value is the number of vertex-disjoint
paths.
\end{algorithmic}
\end{algorithm}

The internal capacity-$1$ edge $v_{\text{in}} \to v_{\text{out}}$ guarantees at
most one unit of flow -- hence at most one path -- passes \emph{through} $v$.
Problem 8 implements this and its tests contrast it with Problem 7: two paths
that are edge-disjoint but share a middle vertex are counted by Problem 7 but not
by Problem 8.

\section{Circulations with Demands and Lower Bounds}

So far every flow had a single source and sink. Many problems instead have
\emph{multiple} supplies and demands and require \emph{minimum} as well as
maximum amounts on edges. The \textbf{circulation} model captures this.

\begin{definition}[Circulation with demands]
Each vertex $v$ has a \emph{demand} $d_v$: if $d_v > 0$, $v$ must \emph{absorb}
$d_v$ units (a sink); if $d_v < 0$, $v$ \emph{supplies} $-d_v$ units (a source).
Each directed edge $e$ has a capacity $c_e$ and possibly a \emph{lower bound}
$\ell_e \ge 0$. A \emph{feasible circulation} assigns a flow $f_e \in [\ell_e,
c_e]$ to every edge so that for each vertex $v$,
\[
   \sum_{e \text{ into } v} f_e \;-\; \sum_{e \text{ out of } v} f_e \;=\; d_v .
\]
A necessary condition is $\sum_v d_v = 0$ (total supply equals total demand).
\end{definition}

\begin{theorem}[Feasibility reduces to max flow]
\label{thm:circulation}
A feasible circulation with demands and lower bounds can be decided by a single
max-flow computation, as follows.
\end{theorem}

\begin{proof}[Construction]
\textbf{Step 1 -- eliminate lower bounds.} For each edge $(u,v)$ with bounds
$[\ell, c]$, ``pre-push'' $\ell$ units: set its capacity to $c - \ell$ and
update demands $d_u \mathrel{+}= \ell$ and $d_v \mathrel{-}= \ell$. (We have
committed $\ell$ units to leave $u$ and arrive at $v$; the residual problem has
lower bound $0$.)

\textbf{Step 2 -- super source and sink.} Add $S$ and $T$. For each vertex $v$
with adjusted demand $d_v < 0$ (net supply), add $S \to v$ with capacity $-d_v$.
For each $v$ with $d_v > 0$ (net demand), add $v \to T$ with capacity $d_v$.

\textbf{Step 3 -- solve.} A feasible circulation exists \emph{iff} the maximum
$S$--$T$ flow \emph{saturates every edge out of $S$} -- equivalently, equals
$D := \sum_{v: d_v > 0} d_v$, the total adjusted demand (and this requires the
adjusted supply and demand to balance). Saturating all source edges means every
supply is fully shipped and every demand fully met, which is exactly a feasible
circulation; conversely a feasible circulation yields such a saturating flow.
\end{proof}

Problem 11 implements this feasibility test, including the lower-bound
elimination and the balance check ($\sum d_v = 0$ after adjustment). This model
is the foundation for many scheduling and routing feasibility questions, and it
generalizes the survey-design and assignment problems of the next section.

\section{Flow Applications I: Bipartite Assignment and Survey Design}

The bipartite-matching reduction generalizes smoothly when each vertex can be
used several times. In \textbf{survey design}, a magazine wants each
\emph{customer} $i$ to be asked between $c_i$ and $c_i'$ questions, and each
\emph{product} $j$ to receive between $p_j$ and $p_j'$ reviews, where customer
$i$ may only be asked about products they have purchased. Model it as a
circulation / bounded flow:
\begin{itemize}
    \item $s \to$ customer $i$ with bounds $[c_i, c_i']$;
    \item customer $i \to$ product $j$ with bounds $[0,1]$ if $i$ purchased $j$;
    \item product $j \to t$ with bounds $[p_j, p_j']$;
    \item a return edge $t \to s$ of unbounded capacity (turning the design into
    a circulation).
\end{itemize}
A feasible integer flow is exactly a valid survey design. This is the
``bounded-degree'' generalization of bipartite matching (Problem 6's assignment
question is the special case where every bound is $[0,1]$ and we ask the flow to
saturate all customers).

\section{Flow Applications II: Baseball Elimination}

A delightful, non-obvious use of max flow decides whether a sports team has been
mathematically \emph{eliminated} from finishing first.

\begin{definition}[Baseball elimination]
We are given teams with current win totals $w_i$ and a schedule of $g_{ij}$
games \emph{remaining} between teams $i$ and $j$. Team $x$ is \emph{eliminated}
if there is \emph{no} way to assign winners to all remaining games so that $x$
finishes with at least as many wins as every other team.
\end{definition}

Let $W = w_x + (\text{remaining games of } x)$ be the most wins $x$ can finish
with. If some team already has $w_i > W$, $x$ is trivially eliminated. Otherwise
build a flow network (Kleinberg \& Tardos, \S 7.12):

\begin{algorithm}[h]
\caption{Baseball Elimination as Max Flow}
\begin{algorithmic}[1]
\For{each unordered pair $\{i,j\}$ of \emph{other} teams with $g_{ij}>0$}
    \State Add a \emph{game node} $\{i,j\}$ with edge $s \to \{i,j\}$ of
    capacity $g_{ij}$.
    \State Add edges $\{i,j\} \to i$ and $\{i,j\} \to j$ of capacity $\infty$.
\EndFor
\For{each other team $i$}
    \State Add edge $i \to t$ of capacity $W - w_i$ \Comment{$i$'s remaining
    ``win budget''}
\EndFor
\State Compute the max $s$--$t$ flow.
\State \Return ``not eliminated'' iff the flow saturates \emph{all} game edges.
\end{algorithmic}
\end{algorithm}

Each game $\{i,j\}$ sends one unit per game played; that unit must flow to the
winner ($i$ or $j$), and the edge $i \to t$ caps how many additional wins team
$i$ may accumulate before passing $W$. If \emph{all} remaining games can be
routed (the flow saturates every $s \to \{i,j\}$ edge), then there is a way to
play out the season with $x$ tied for first or better; otherwise $x$ is
eliminated. Problem 9 implements this on a small dataset, including the trivial
elimination shortcut.

The minimum cut, by the way, yields a \emph{certificate of elimination}: a
subset $R$ of teams whose combined current and remaining intra-group wins force
their average above $W$ -- a human-readable reason why $x$ cannot catch up.

\section{Flow Applications III: Project Selection (Max-Weight Closure)}

\begin{definition}[Project selection]
We have projects, each with a profit $p_v$ (possibly negative -- a cost).
Some projects have \emph{prerequisites}: an edge $(p,q)$ means ``selecting $p$
requires also selecting $q$.'' A set $S$ of projects is \emph{feasible} if it is
\emph{closed} under prerequisites (if $p \in S$ and $(p,q)$ is a prerequisite,
then $q \in S$). The goal is a feasible $S$ maximizing $\sum_{v \in S} p_v$.
\end{definition}

This is the \textbf{maximum-weight closure} problem, and it is solved by a
\emph{minimum cut} (Kleinberg \& Tardos, \S 7.11):

\begin{algorithm}[h]
\caption{Project Selection via Min-Cut}
\begin{algorithmic}[1]
\State Let $P = \sum_{v : p_v > 0} p_v$ (sum of positive profits).
\For{each project $v$ with $p_v > 0$} add $s \to v$ with capacity $p_v$.
\EndFor
\For{each project $v$ with $p_v < 0$} add $v \to t$ with capacity $-p_v$.
\EndFor
\For{each prerequisite $(p,q)$} add $p \to q$ with capacity $\infty$.
\EndFor
\State Compute the minimum $s$--$t$ cut, of value $C$.
\State \Return profit $P - C$, and $S = \{v : v \text{ on the source side}\}$.
\end{algorithmic}
\end{algorithm}

\begin{theorem}
\label{thm:closure}
The maximum profit equals $P - C$, where $C$ is the minimum cut, and the
source side of the minimum cut is an optimal feasible project set.
\end{theorem}

\begin{proof}[Sketch]
The infinite-capacity prerequisite edges guarantee that any \emph{finite} cut
never separates a selected project from a prerequisite -- so the source side
$S$ is always closed. For such a cut, the cut capacity equals
$\sum_{v \in S, p_v < 0}(-p_v) + \sum_{v \notin S, p_v > 0} p_v$, i.e.\ the
forgone positive profits plus the incurred costs. Minimizing this is the same as
maximizing $\sum_{v \in S} p_v = P - (\text{cut})$. Hence min-cut $\Leftrightarrow$
max-profit closure.
\end{proof}

Problem 10 implements both the min-cut solution and a brute-force search over all
closed subsets, and the tests confirm they agree on random small instances. The
source side of the cut is recovered with \texttt{MaxFlow.min\_cut\_reachable}.

\section{Flow Applications IV: Image Segmentation as Min-Cut}

A final, visually compelling application: \textbf{image segmentation}. We wish to
label each pixel as \emph{foreground} or \emph{background}. Each pixel $i$ has a
\emph{likelihood} $a_i \ge 0$ of being foreground and $b_i \ge 0$ of being
background, and each pair of neighboring pixels $(i,j)$ has a
\emph{separation penalty} $p_{ij} \ge 0$ we pay if we split them across the two
labels (this enforces smooth, contiguous regions).

We \emph{maximize} $\sum_{i \in F} a_i + \sum_{j \in B} b_j - \sum_{(i,j) \text{
split}} p_{ij}$, which is equivalent to \emph{minimizing}
$\sum_{i \in F} b_i + \sum_{j \in B} a_j + \sum_{(i,j) \text{ split}} p_{ij}$.
Build a network with $s$ = ``foreground terminal,'' $t$ = ``background
terminal'':
\begin{itemize}
    \item $s \to i$ of capacity $a_i$ for every pixel $i$;
    \item $i \to t$ of capacity $b_i$ for every pixel $i$;
    \item between neighbors, \emph{two} directed edges $i \to j$ and $j \to i$,
    each of capacity $p_{ij}$.
\end{itemize}
A minimum $s$--$t$ cut places foreground pixels on the source side and background
on the sink side; its capacity is exactly the quantity to be minimized, with the
neighbor edges contributing $p_{ij}$ precisely when $i$ and $j$ end up on
opposite sides. Thus image segmentation is \emph{another} min-cut problem -- the
same machinery as project selection, with a different network. (The shared
\texttt{MaxFlow.min\_cut\_reachable} helper would recover the foreground set.)

\section{Summary and Looking Ahead}

This week we treated max flow as a \textbf{modeling and reduction tool}, and saw
how a single black-box algorithm (Edmonds--Karp, from Week 10) solves a wide
range of combinatorial problems:

\begin{itemize}
    \item \textbf{Bipartite matching} $=$ unit-capacity max flow
    (Theorem~\ref{thm:matching-flow}); Problems 2--3, 6.
    \item \textbf{Hall's theorem} (Theorem~\ref{thm:hall}) and \textbf{K\H{o}nig's
    theorem} (Theorem~\ref{thm:konig}): perfect matchings via the deficiency
    condition, and $\max$-matching $=$ $\min$-vertex-cover in bipartite graphs;
    Problems 4, 5, 12.
    \item \textbf{Disjoint paths} (Menger, Theorem~\ref{thm:menger-edge}):
    edge-disjoint via unit capacities, vertex-disjoint via node splitting;
    Problems 7--8.
    \item \textbf{Circulations with demands and lower bounds}
    (Theorem~\ref{thm:circulation}): feasibility by one max-flow call; Problem 11.
    \item \textbf{Min-cut applications}: baseball elimination (Problem 9),
    project selection / max-weight closure (Theorem~\ref{thm:closure}, Problem
    10), and image segmentation -- all decided by a minimum cut.
\end{itemize}

The unifying lesson is that the \emph{hard work is the reduction}: encoding
constraints as capacities, recognizing when a min-max theorem applies, and
proving the correspondence between combinatorial objects and flows/cuts. Once
that is done, optimality and a polynomial running time come for free from the
max-flow min-cut theorem.

\textbf{Looking ahead.} Flow has given us a large catalog of problems that are
\emph{efficiently solvable}. Next week we turn to the other side of the coin:
\textbf{intractability} and \textbf{NP-completeness} (Kleinberg \& Tardos, Ch.\
8). We will study \emph{reductions} again -- but now to argue that certain
problems are \emph{at least as hard} as a whole class of others, so that a
polynomial algorithm for one would solve them all. Many problems that look
deceptively close to the matchings and cuts of this week (e.g.\ maximum matching
in \emph{general} graphs is still polynomial, but maximum independent set, or
3-dimensional matching, are NP-complete) sit just beyond the reach of flow --
and understanding exactly where that boundary lies is the subject of the rest of
the course.

\end{document}
